%%% FIELD proof-zero-gamma-reduction
Assume \(\Gamma=0\). Fix \(S\in\mathcal P_k\), \(\pi\in\mathcal F\), \(x\in\mathcal X\), and \(c\in\mathcal C\). By \(\mathrm{def{:}score\text{-}cost}\), if \(c(x)\neq\pi(x)\), then \(d_{\pi xc}=1\). If \(c(x)=\pi(x)=a_0\), the same definition gives \(d_{\pi xc}=0\). In the remaining case, \(c(x)=\pi(x)\neq a_0\), and therefore
\[
d_{\pi xc}=\min\{1,\Gamma|m_{\pi}-m_c|\}=\min\{1,0\}=0.
\]
Consequently,
\[
d_{\pi xc}=
\begin{cases}
0,&c(x)=\pi(x),\\
1,&c(x)\neq\pi(x).
\end{cases}
\]
By the definition of \(r_{\pi x}(S)\) in \(\mathrm{def{:}score\text{-}cost}\),
\[
r_{\pi x}(S)=\min\bigl(\{1\}\cup\{d_{\pi xc}:c\in S\}\bigr).
\]
If some \(c\in S\) satisfies \(c(x)=\pi(x)\), the set inside the minimum contains zero, so \(r_{\pi x}(S)=0\). If no such \(c\) exists, including when \(S=\varnothing\), every displayed cost contributed by \(S\) equals one and the adjoined element is one, so \(r_{\pi x}(S)=1\). Hence
\[
r_{\pi x}(S)=\mathbf 1\{c(x)\neq\pi(x)\text{ for every }c\in S\}.
\]
The already-established width identity \(\mathrm{thm{:}width\text{-}identity}\) now yields
\[
\overline W(S)
=\max_{\pi\in\mathcal F}R_{\pi}(S)
=\max_{\pi\in\mathcal F}\sum_{x\in\mathcal X}q_x r_{\pi x}(S)
=\max_{\pi\in\mathcal F}\sum_{x\in\mathcal X}q_x\mathbf 1\{c(x)\neq\pi(x)\text{ for every }c\in S\}.
\]
Finally, suppose that every future cell is exactly covered: for each \(\pi\in\mathcal F\) and \(x\in\mathcal X\), there exists \(c\in S\) with \(c(x)=\pi(x)\). Then every indicator in the last display is zero. Thus each weighted sum is zero, and because \(\mathcal F\) is finite and nonempty, its maximum is zero. Therefore \(\overline W(S)=0\), as claimed.
